Capítulo  2
 Organização  do  sistema  operacional
 Um  requisito  fundamental  para  um  sistema  operacional  é  suportar  diversas  atividades  simultaneamente.  Por  exemplo,  
usando  a  interface  de  chamada  de  sistema  descrita  no  Capítulo  1,  um  processo  pode  iniciar  novos  processos  com  fork.
 O  sistema  operacional  deve  compartilhar  o  tempo  dos  recursos  do  computador  entre  esses  processos.  Por  exemplo,  
mesmo  que  haja  mais  processos  do  que  CPUs  de  hardware,  o  sistema  operacional  deve  garantir  que  todos  os  processos  
tenham  a  chance  de  executar.  O  sistema  operacional  também  deve  providenciar  o  isolamento  entre  os  processos.  Ou  
seja,  se  um  processo  tiver  um  bug  e  apresentar  mau  funcionamento,  isso  não  deve  afetar  os  processos  que  não  dependem  
do  processo  com  bug.  O  isolamento  completo,  no  entanto,  é  muito  forte,  pois  deve  ser  possível  que  os  processos  interajam  
intencionalmente;  pipelines  são  um  exemplo.  Portanto,  um  sistema  operacional  deve  atender  a  três  requisitos:  
multiplexação,  isolamento  e  interação.
 Este  capítulo  fornece  uma  visão  geral  de  como  os  sistemas  operacionais  são  organizados  para  atender  a  esses  
três  requisitos.  Existem  muitas  maneiras  de  fazer  isso,  mas  este  texto  se  concentra  em  projetos  tradicionais  centrados  
em  um  kernel  monolítico,  usado  por  muitos  sistemas  operacionais  Unix.  Este  capítulo  também  fornece  uma  visão  geral  
de  um  processo  xv6,  que  é  a  unidade  de  isolamento  no  xv6,  e  da  criação  do  primeiro  processo  quando  o  xv6  é  iniciado.
 O  Xv6  é  executado  em  um  microprocessador  RISC-V  multi-core1 ,  e  grande  parte  de  sua  funcionalidade  de  baixo  
nível  (por  exemplo,  sua  implementação  de  processo)  é  específica  para  RISC-V.  RISC-V  é  uma  CPU  de  64  bits,  e  o  xv6  
é  escrito  em  C  "LP64",  o  que  significa  que  long  (L)  e  ponteiros  (P)  na  linguagem  de  programação  C  são  de  64  bits,  mas  
int  é  de  32  bits.  Este  livro  pressupõe  que  o  leitor  tenha  feito  um  pouco  de  programação  em  nível  de  máquina  em  
alguma  arquitetura,  e  apresentará  ideias  específicas  do  RISC-V  à  medida  que  surgirem.  Uma  referência  útil  para  RISC
V  é  "The  RISC-V  Reader:  An  Open  Architecture  Atlas"  [15].  O  ISA  em  nível  de  usuário  [2]  e  a  arquitetura  privilegiada  
[3]  são  as  especificações  oficiais.
 A  CPU  de  um  computador  completo  é  cercada  por  hardware  de  suporte,  grande  parte  dele  na  forma  de  interfaces  
de  E/S.  O  Xv6  foi  escrito  para  o  hardware  de  suporte  simulado  pela  opção  "-machine  virt"  do  qemu.  Isso  inclui  RAM,  
uma  ROM  contendo  o  código  de  inicialização,  uma  conexão  serial  com  o  teclado/tela  do  usuário  e  um  disco  para  
armazenamento.
 1Por  "multinúcleo",  este  texto  se  refere  a  múltiplas  CPUs  que  compartilham  memória,  mas  executam  em  paralelo,  cada  uma  com  seu  próprio  
conjunto  de  registradores.  Este  texto  às  vezes  usa  o  termo  multiprocessador  como  sinônimo  de  multinúcleo,  embora  multiprocessador  também  
possa  se  referir  mais  especificamente  a  um  computador  com  vários  chips  de  processador  distintos.
 21
Machine Translated by Google
 2.1  Abstraindo  recursos  físicos
 A  primeira  pergunta  que  se  pode  fazer  ao  se  deparar  com  um  sistema  operacional  é:  por  que  tê-lo?  Ou  seja,  seria  possível  
implementar  as  chamadas  de  sistema  da  Figura  1.2  como  uma  biblioteca,  com  a  qual  os  aplicativos  se  conectam.  Nesse  plano,  
cada  aplicativo  poderia  até  ter  sua  própria  biblioteca,  adaptada  às  suas  necessidades.  Os  aplicativos  poderiam  interagir  
diretamente  com  os  recursos  de  hardware  e  utilizá-los  da  melhor  maneira  possível  para  o  aplicativo  (por  exemplo,  para  atingir  um  
desempenho  alto  ou  previsível).  Alguns  sistemas  operacionais  para  dispositivos  embarcados  ou  sistemas  de  tempo  real  são  
organizados  dessa  maneira.
 A  desvantagem  dessa  abordagem  de  biblioteca  é  que,  se  houver  mais  de  um  aplicativo  em  execução,  eles  precisam  se  
comportar  bem.  Por  exemplo,  cada  aplicativo  precisa  liberar  a  CPU  periodicamente  para  que  outros  aplicativos  possam  ser  
executados.  Esse  esquema  cooperativo  de  compartilhamento  de  tempo  pode  ser  aceitável  se  todos  os  aplicativos  confiarem  uns  
nos  outros  e  não  apresentarem  bugs.  É  mais  comum  que  aplicativos  não  confiem  uns  nos  outros  e  apresentem  bugs,  portanto,  
muitas  vezes,  deseja-se  um  isolamento  mais  forte  do  que  o  proporcionado  por  um  esquema  cooperativo.
 Para  obter  um  isolamento  forte,  é  útil  proibir  que  aplicativos  acessem  diretamente  recursos  de  hardware  sensíveis  e,  
em  vez  disso,  abstrair  os  recursos  em  serviços.  Por  exemplo,  aplicativos  Unix  interagem  com  o  armazenamento  apenas  
por  meio  das  chamadas  de  sistema  de  abertura,  leitura,  gravação  e  fechamento  do  sistema  de  arquivos,  em  vez  de  ler  e  
gravar  diretamente  no  disco.  Isso  proporciona  ao  aplicativo  a  conveniência  de  nomes  de  caminho  e  permite  que  o  
sistema  operacional  (como  implementador  da  interface)  gerencie  o  disco.  Mesmo  que  o  isolamento  não  seja  uma  
preocupação,  programas  que  interagem  intencionalmente  (ou  apenas  desejam  se  manter  fora  do  caminho  uns  dos  outros)  
provavelmente  considerarão  um  sistema  de  arquivos  uma  abstração  mais  conveniente  do  que  o  uso  direto  do  disco.
 Da  mesma  forma,  o  Unix  alterna  de  forma  transparente  as  CPUs  de  hardware  entre  os  processos,  salvando  e  restaurando  
o  estado  dos  registradores  conforme  necessário,  para  que  os  aplicativos  não  precisem  se  preocupar  com  o  compartilhamento  de  
tempo.  Essa  transparência  permite  que  o  sistema  operacional  compartilhe  CPUs  mesmo  que  alguns  aplicativos  estejam  em  loops  
infinitos.
 Como  outro  exemplo,  os  processos  Unix  usam  o  comando  exec  para  construir  sua  imagem  de  memória,  em  vez  de  interagir  
diretamente  com  a  memória  física.  Isso  permite  que  o  sistema  operacional  decida  onde  colocar  um  processo  na  memória;  se  a  
memória  estiver  limitada,  o  sistema  operacional  pode  até  armazenar  alguns  dados  de  um  processo  em  disco.  O  comando  exec  
também  oferece  aos  usuários  a  conveniência  de  um  sistema  de  arquivos  para  armazenar  imagens  de  programas  executáveis.
 Muitas  formas  de  interação  entre  processos  Unix  ocorrem  por  meio  de  descritores  de  arquivo.  Os  descritores  de  arquivo  não  
apenas  abstraem  muitos  detalhes  (por  exemplo,  onde  os  dados  em  um  pipe  ou  arquivo  são  armazenados),  como  também  são  
definidos  de  forma  a  simplificar  a  interação.  Por  exemplo,  se  uma  aplicação  em  um  pipeline  falhar,  o  kernel  gera  um  sinal  de  fim  
de  arquivo  para  o  próximo  processo  no  pipeline.
 A  interface  de  chamada  de  sistema  na  Figura  1.2  foi  cuidadosamente  projetada  para  proporcionar  conveniência  ao  
programador  e  a  possibilidade  de  forte  isolamento.  A  interface  Unix  não  é  a  única  maneira  de  abstrair  recursos,  mas  provou  ser  
uma  ótima  maneira.
 2.2  Modo  de  usuário,  modo  supervisor  e  chamadas  de  sistema
 O  isolamento  forte  requer  uma  fronteira  rígida  entre  os  aplicativos  e  o  sistema  operacional.  Se  o  aplicativo  cometer  um  erro,  não  
queremos  que  o  sistema  operacional  falhe  ou  que  outros  aplicativos  sejam  afetados.
 22
Machine Translated by Google
 falha.  Em  vez  disso,  o  sistema  operacional  deve  ser  capaz  de  limpar  o  aplicativo  com  falha  e  continuar  executando  
outros  aplicativos.  Para  alcançar  um  isolamento  forte,  o  sistema  operacional  deve  garantir  que  os  aplicativos  não  
possam  modificar  (ou  mesmo  ler)  suas  estruturas  de  dados  e  instruções,  e  que  os  aplicativos  não  possam  acessar  a  
memória  de  outros  processos.
 As  CPUs  fornecem  suporte  de  hardware  para  isolamento  robusto.  Por  exemplo,  o  RISC-V  possui  três  modos  nos  
quais  a  CPU  pode  executar  instruções:  modo  máquina,  modo  supervisor  e  modo  usuário.  Instruções  de  operação  
executadas  em  modo  máquina  têm  privilégio  total;  uma  CPU  inicia  em  modo  máquina.  O  modo  máquina  destina-se  
principalmente  à  configuração  de  um  computador.  O  Xv6  executa  algumas  linhas  em  modo  máquina  e  depois  muda  
para  o  modo  supervisor.
 No  modo  supervisor,  a  CPU  tem  permissão  para  executar  instruções  privilegiadas:  por  exemplo,  habilitar  e  
desabilitar  interrupções,  ler  e  escrever  no  registrador  que  contém  o  endereço  de  uma  tabela  de  páginas,  etc.  Se  um  
aplicativo  no  modo  usuário  tentar  executar  uma  instrução  privilegiada,  a  CPU  não  executa  a  instrução,  mas  alterna  
para  o  modo  supervisor  para  que  o  código  do  modo  supervisor  possa  encerrar  o  aplicativo,  porque  ele  fez  algo  que  não  
deveria.  A  Figura  1.1  no  Capítulo  1  ilustra  essa  organização.  Um  aplicativo  pode  executar  apenas  instruções  do  modo  
usuário  (por  exemplo,  somar  números,  etc.)  e  é  dito  estar  em  execução  no  espaço  do  usuário,  enquanto  o  software  no  
modo  supervisor  também  pode  executar  instruções  privilegiadas  e  é  dito  estar  em  execução  no  espaço  do  kernel.  O  
software  em  execução  no  espaço  do  kernel  (ou  no  modo  supervisor)  é  chamado  de  kernel.
 Um  aplicativo  que  deseja  invocar  uma  função  do  kernel  (por  exemplo,  a  chamada  de  sistema  read  no  xv6)  deve  
fazer  a  transição  para  o  kernel;  um  aplicativo  não  pode  invocar  uma  função  do  kernel  diretamente.  As  CPUs  fornecem  
uma  instrução  especial  que  alterna  a  CPU  do  modo  usuário  para  o  modo  supervisor  e  entra  no  kernel  em  um  ponto  de  
entrada  especificado  pelo  kernel.  (O  RISC-V  fornece  a  instrução  ecall  para  esse  propósito.)
 Após  a  CPU  alternar  para  o  modo  supervisor,  o  kernel  pode  validar  os  argumentos  da  chamada  de  sistema  (por  
exemplo,  verificar  se  o  endereço  passado  para  a  chamada  de  sistema  faz  parte  da  memória  do  aplicativo),  decidir  se  o  
aplicativo  tem  permissão  para  executar  a  operação  solicitada  (por  exemplo,  verificar  se  o  aplicativo  tem  permissão  para  
gravar  o  arquivo  especificado)  e,  em  seguida,  negá-la  ou  executá-la.  É  importante  que  o  kernel  controle  o  ponto  de  
entrada  para  transições  para  o  modo  supervisor;  se  o  aplicativo  pudesse  decidir  o  ponto  de  entrada  do  kernel,  um  
aplicativo  malicioso  poderia,  por  exemplo,  entrar  no  kernel  em  um  ponto  onde  a  validação  dos  argumentos  é  ignorada.
 2.3  Organização  do  kernel
 Uma  questão  fundamental  de  projeto  é  qual  parte  do  sistema  operacional  deve  ser  executada  no  modo  supervisor.  
Uma  possibilidade  é  que  todo  o  sistema  operacional  resida  no  kernel,  de  modo  que  as  implementações  de  todas  as  
chamadas  de  sistema  sejam  executadas  no  modo  supervisor.  Essa  organização  é  chamada  de  kernel  monolítico.
 Nesta  organização,  todo  o  sistema  operacional  é  executado  com  privilégios  totais  de  hardware.  Essa  organização  
é  conveniente  porque  o  projetista  do  sistema  operacional  não  precisa  decidir  qual  parte  do  sistema  operacional  não  
precisa  de  privilégios  totais  de  hardware.  Além  disso,  facilita  a  cooperação  entre  diferentes  partes  do  sistema  
operacional.  Por  exemplo,  um  sistema  operacional  pode  ter  um  cache  de  buffer  que  pode  ser  compartilhado  tanto  pelo  
sistema  de  arquivos  quanto  pelo  sistema  de  memória  virtual.
 Uma  desvantagem  da  organização  monolítica  é  que  as  interfaces  entre  as  diferentes  partes  do  sistema  operacional  
são  frequentemente  complexas  (como  veremos  no  restante  deste  texto)  e,  portanto,  é
 23
Machine Translated by Google
 usuário
 espaço
 núcleo
 espaço
 concha
 Servidor  de  arquivos
 Enviar  mensagem
 Microkernel
 Figura  2.1:  Um  microkernel  com  um  servidor  de  sistema  de  arquivos
 É  fácil  para  um  desenvolvedor  de  sistema  operacional  cometer  um  erro.  Em  um  kernel  monolítico,  um  erro  é  fatal,  pois  um  erro  
no  modo  supervisor  frequentemente  causa  falha  no  kernel.  Se  o  kernel  falhar,  o  computador  para  de  funcionar  e,  portanto,  todos  
os  aplicativos  também  falham.  O  computador  precisa  ser  reiniciado  para  iniciar  novamente.
 Para  reduzir  o  risco  de  erros  no  kernel,  os  projetistas  de  sistemas  operacionais  podem  minimizar  a  quantidade  de  código  do  
sistema  operacional  executado  no  modo  supervisor  e  executar  a  maior  parte  do  sistema  operacional  no  modo  usuário.  Essa  
organização  do  kernel  é  chamada  de  microkernel.
 A  Figura  2.1  ilustra  este  projeto  de  microkernel.  Na  figura,  o  sistema  de  arquivos  é  executado  como  um  processo  em  nível  
de  usuário.  Os  serviços  do  sistema  operacional  executados  como  processos  são  chamados  de  servidores.  Para  permitir  que  os  
aplicativos  interajam  com  o  servidor  de  arquivos,  o  kernel  fornece  um  mecanismo  de  comunicação  entre  processos  para  enviar  
mensagens  de  um  processo  em  modo  de  usuário  para  outro.  Por  exemplo,  se  um  aplicativo,  como  o  shell,  deseja  ler  ou  gravar  
um  arquivo,  ele  envia  uma  mensagem  ao  servidor  de  arquivos  e  aguarda  uma  resposta.
 Em  um  microkernel,  a  interface  do  kernel  consiste  em  algumas  funções  de  baixo  nível  para  iniciar  aplicativos,  enviar  
mensagens,  acessar  hardware  do  dispositivo,  etc.  Essa  organização  permite  que  o  kernel  seja  relativamente  simples,  já  que  a  
maior  parte  do  sistema  operacional  reside  em  servidores  de  nível  de  usuário.
 No  mundo  real,  tanto  kernels  monolíticos  quanto  microkernels  são  populares.  Muitos  kernels  Unix  são  monolíticos.  Por  
exemplo,  o  Linux  possui  um  kernel  monolítico,  embora  algumas  funções  do  sistema  operacional  sejam  executadas  como  
servidores  em  nível  de  usuário  (por  exemplo,  o  sistema  de  janelas).  O  Linux  oferece  alto  desempenho  para  aplicativos  com  uso  
intensivo  de  sistema  operacional,  em  parte  porque  os  subsistemas  do  kernel  podem  ser  totalmente  integrados.
 Sistemas  operacionais  como  Minix,  L4  e  QNX  são  organizados  como  um  microkernel  com  servidores  e  têm  sido  amplamente  
implantados  em  ambientes  embarcados.  Uma  variante  do  L4,  seL4,  é  pequena  o  suficiente  para  ter  sua  segurança  de  memória  e  
outras  propriedades  de  segurança  verificadas  [8].
 Há  muito  debate  entre  desenvolvedores  de  sistemas  operacionais  sobre  qual  organização  é  melhor,  e  não  há  evidências  
conclusivas  de  um  lado  ou  de  outro.  Além  disso,  depende  muito  do  que  "melhor"  significa:  desempenho  mais  rápido,  código  
menor,  confiabilidade  do  kernel,  confiabilidade  de  todo  o  sistema  operacional  (incluindo  serviços  em  nível  de  usuário),  etc.
 Há  também  considerações  práticas  que  podem  ser  mais  importantes  do  que  a  questão  de  qual  organização.  Alguns  sistemas  
operacionais  possuem  um  microkernel,  mas  executam  alguns  dos  serviços  de  nível  de  usuário  no  espaço  do  kernel  por  questões  
de  desempenho.  Alguns  sistemas  operacionais  possuem  kernels  monolíticos  porque  foi  assim  que  começaram,  e  há  pouco  
incentivo  para  migrar  para  uma  organização  puramente  de  microkernel,  pois  novos  recursos  podem  ser  mais  importantes  do  que  
reescrever  o  sistema  operacional  existente  para  se  adequar  a  um  design  de  microkernel.
 Da  perspectiva  deste  livro,  microkernel  e  sistemas  operacionais  monolíticos  compartilham  muitas  ideias-chave.  Eles  
implementam  chamadas  de  sistema,  usam  tabelas  de  páginas,  lidam  com  interrupções  e  oferecem  suporte
 24
Machine Translated by Google
 Arquivo
 bio.c  
console.c  
entrada.S
 exec.c
 arquivo.c  fs.c
 Descrição
 Cache  de  bloco  de  disco  para  o  sistema  de  arquivos.
 Conecte-se  ao  teclado  e  à  tela  do  usuário.
 Primeiras  instruções  de  inicialização.
 chamada  de  sistema  exec().
 Suporte  ao  descritor  de  arquivo.
 Sistema  de  arquivos.
 kalloc.c  Alocador  de  páginas  físicas.
 kernelvec.S  Lida  com  armadilhas  do  kernel  e  interrupções  de  temporizador.
 log.c  Registro  do  sistema  de  arquivos  e  recuperação  de  falhas.
 main.c  Controla  a  inicialização  de  outros  módulos  durante  a  inicialização.
 pipe.c  Tubos.
 plic.c  Controlador  de  interrupção  RISC-V.
 printf.c  Saída  formatada  para  o  console.
 Processos  e  agendamento.
 proc.c  
sleeplock.c  Bloqueios  que  produzem  a  CPU.
 spinlock.c  Bloqueios  que  não  cedem  a  CPU.
 start.c  
string.c  
swtch.S  
syscall.c  
sysfile.c
 Código  de  inicialização  inicial  do  modo  máquina.
 Biblioteca  de  strings  e  matrizes  de  bytes  em  C.
 Troca  de  thread.
 Despachar  chamadas  do  sistema  para  a  função  de  manuseio.
 Chamadas  de  sistema  relacionadas  a  arquivos.
 Chamadas  de  sistema  relacionadas  ao  processo.
 sysproc.c  
trampoline.S  Código  de  montagem  para  alternar  entre  usuário  e  kernel.
 trap.c  
uart.c
 Código  C  para  manipular  e  retornar  armadilhas  e  interrupções.
 Driver  de  dispositivo  de  console  de  porta  serial.
 virtio_disk.c  Driver  de  dispositivo  de  disco.
 vm.c
 Gerenciar  tabelas  de  páginas  e  espaços  de  endereço.
 Figura  2.2:  Arquivos  de  origem  do  kernel  Xv6.
 processos,  eles  usam  bloqueios  para  controle  de  simultaneidade,  eles  implementam  um  sistema  de  arquivos,  etc.  Este  livro
 concentra-se  nessas  ideias  principais.
 O  Xv6  é  implementado  como  um  kernel  monolítico,  como  a  maioria  dos  sistemas  operacionais  Unix.  Assim,  o  xv6
 A  interface  do  kernel  corresponde  à  interface  do  sistema  operacional,  e  o  kernel  implementa  o  sistema  operacional  completo.  Como  o  xv6  
não  fornece  muitos  serviços,  seu  kernel  é  menor  do  que  alguns
 microkernels,  mas  conceitualmente  o  xv6  é  monolítico.
 2.4  Código:  organização  xv6
 O  código-fonte  do  kernel  xv6  está  no  subdiretório  kernel/.  O  código-fonte  é  dividido  em  arquivos,  a  seguir
 uma  noção  aproximada  de  modularidade;  a  Figura  2.2  lista  os  arquivos.  As  interfaces  entre  módulos  são  definidas  em
 25
Machine Translated by Google
 defs.h  (kernel/defs.h).
 2.5  Visão  geral  do  processo
 MAXVA
 0
 estrutura  de  
armadilha  de  trampolim
 monte
 pilha  de  usuário
 texto  do  usuário
 e  dados
 Figura  2.3:  Layout  do  espaço  de  endereço  virtual  de  um  processo
 A  unidade  de  isolamento  no  xv6  (assim  como  em  outros  sistemas  operacionais  Unix)  é  um  processo.  A  abstração  de  
processo  impede  que  um  processo  destrua  ou  espione  a  memória,  a  CPU,  os  descritores  de  arquivo,  etc.  de  outro  processo.  
Também  impede  que  um  processo  destrua  o  próprio  kernel,  impedindo  que  um  processo  subverta  os  mecanismos  de  
isolamento  do  kernel.  O  kernel  deve  implementar  a  abstração  de  processo  com  cuidado,  pois  um  aplicativo  malicioso  ou  
com  bugs  pode  induzir  o  kernel  ou  o  hardware  a  fazer  algo  errado  (por  exemplo,  burlar  o  isolamento).  Os  mecanismos  
usados  pelo  kernel  para  implementar  processos  incluem  o  sinalizador  de  modo  usuário/supervisor,  espaços  de  endereço  e  
fatiamento  de  tempo  de  threads.
 Para  ajudar  a  reforçar  o  isolamento,  a  abstração  do  processo  dá  a  um  programa  a  ilusão  de  que  ele  possui  sua  própria  
máquina  privada.  Um  processo  fornece  ao  programa  o  que  parece  ser  um  sistema  de  memória  privado,  ou  espaço  de  
endereço,  que  outros  processos  não  conseguem  ler  ou  escrever.  Um  processo  também  fornece  ao  programa  o  que  parece  
ser  sua  própria  CPU  para  executar  as  instruções  do  programa.
 O  Xv6  utiliza  tabelas  de  páginas  (implementadas  por  hardware)  para  dar  a  cada  processo  seu  próprio  espaço  de  
endereço.  A  tabela  de  páginas  RISC-V  traduz  (ou  "mapeia")  um  endereço  virtual  (o  endereço  que  uma  instrução  RISC-V  
manipula)  para  um  endereço  físico  (um  endereço  que  o  chip  da  CPU  envia  para  a  memória  principal).
 O  Xv6  mantém  uma  tabela  de  páginas  separada  para  cada  processo  que  define  o  espaço  de  endereço  desse  processo.
 Conforme  ilustrado  na  Figura  2.3,  um  espaço  de  endereço  inclui  a  memória  do  usuário  do  processo,  começando  no  endereço  
virtual  zero.  As  instruções  vêm  primeiro,  seguidas  pelas  variáveis  globais,  depois  pela  pilha  e,  finalmente,  por  uma  área  de  
"heap"  (para  malloc)  que  o  processo  pode  expandir  conforme  necessário.  Há  uma  série  de  fatores  que  limitam  o  tamanho  
máximo  do  espaço  de  endereço  de  um  processo:  ponteiros  no  RISC-V  têm  64  bits  de  largura;  o  hardware  usa  apenas  os  39  
bits  mais  baixos  ao  procurar  endereços  virtuais  em  tabelas  de  páginas;  e  o  xv6  usa  apenas  38  desses  39  bits.  Portanto,  o  
endereço  máximo  é  2  ÿ  1  =  0x3ffffffffff,  que  é
 38
 26
Machine Translated by Google
 MAXVA  (kernel/riscv.h:363).  No  topo  do  espaço  de  endereço,  o  xv6  reserva  uma  página  para  um  trampolim  e  uma  
página  que  mapeia  o  trapframe  do  processo.  O  xv6  usa  essas  duas  páginas  para  fazer  a  transição  de  entrada  e  
saída  do  kernel;  a  página  do  trampolim  contém  o  código  para  a  transição  de  entrada  e  saída  do  kernel,  e  o  
mapeamento  do  trapframe  é  necessário  para  salvar/restaurar  o  estado  do  processo  do  usuário,  como  explicaremos  no  Capítulo  4.
 O  kernel  xv6  mantém  muitas  partes  de  estado  para  cada  processo,  que  ele  reúne  em  uma  estrutura  proc
 (kernel/proc.h:85).  As  partes  mais  importantes  do  estado  do  kernel  de  um  processo  são  sua  tabela  de  páginas,  sua  pilha  
do  kernel  e  seu  estado  de  execução.  Usaremos  a  notação  p->xxx  para  nos  referirmos  aos  elementos  da  estrutura  proc ;  
por  exemplo,  p->pagetable  é  um  ponteiro  para  a  tabela  de  páginas  do  processo.
 Cada  processo  possui  uma  thread  de  execução  (ou  thread,  para  abreviar)  que  executa  as  instruções  do  processo.  
Uma  thread  pode  ser  suspensa  e  posteriormente  retomada.  Para  alternar  entre  processos  de  forma  transparente,  o  kernel  
suspende  a  thread  em  execução  e  retoma  a  thread  de  outro  processo.  Grande  parte  do  estado  de  uma  thread  (variáveis  
locais,  endereços  de  retorno  de  chamadas  de  função)  é  armazenada  nas  pilhas  da  thread.
 Cada  processo  possui  duas  pilhas:  uma  pilha  do  usuário  e  uma  pilha  do  kernel  (p->kstack).  Quando  o  processo  está  
executando  instruções  do  usuário,  apenas  sua  pilha  do  usuário  está  em  uso,  e  sua  pilha  do  kernel  está  vazia.  Quando  o  
processo  entra  no  kernel  (para  uma  chamada  de  sistema  ou  interrupção),  o  código  do  kernel  é  executado  na  pilha  do  kernel  
do  processo;  enquanto  um  processo  está  no  kernel,  sua  pilha  do  usuário  ainda  contém  dados  salvos,  mas  não  é  usada  
ativamente.  A  thread  de  um  processo  alterna  entre  usar  ativamente  sua  pilha  do  usuário  e  sua  pilha  do  kernel.
 A  pilha  do  kernel  é  separada  (e  protegida  do  código  do  usuário)  para  que  o  kernel  possa  ser  executado  mesmo  que  um  
processo  tenha  destruído  sua  pilha  do  usuário.
 Um  processo  pode  fazer  uma  chamada  de  sistema  executando  a  instrução  ecall  do  RISC-V .  Essa  instrução  aumenta  
o  nível  de  privilégio  de  hardware  e  altera  o  contador  de  programa  para  um  ponto  de  entrada  definido  pelo  kernel.
 O  código  no  ponto  de  entrada  alterna  para  uma  pilha  do  kernel  e  executa  as  instruções  do  kernel  que  implementam  a  
chamada  do  sistema.  Quando  a  chamada  do  sistema  é  concluída,  o  kernel  retorna  para  a  pilha  do  usuário  e  retorna  ao  
espaço  do  usuário  chamando  a  instrução  sret ,  que  reduz  o  nível  de  privilégio  de  hardware  e  retoma  a  execução  das  
instruções  do  usuário  logo  após  a  instrução  de  chamada  do  sistema.  A  thread  de  um  processo  pode  "bloquear"  no  kernel  
para  aguardar  a  E/S  e  retomar  de  onde  parou  quando  a  E/S  terminar.
 p->state  indica  se  o  processo  está  alocado,  pronto  para  execução,  em  execução,  aguardando  E/S  ou  encerrando.  
p
>pagetable  contém  a  tabela  de  páginas  do  processo,  no  formato  esperado  pelo  hardware  RISC-V.  O  Xv6  faz  com  
que  o  hardware  de  paginação  utilize  a  p->pagetable  de  um  processo  ao  executá-lo  no  espaço  do  usuário.  A  tabela  
de  páginas  de  um  processo  também  serve  como  registro  dos  endereços  das  páginas  físicas  alocadas  para  
armazenar  a  memória  do  processo.
 Em  resumo,  um  processo  agrupa  duas  ideias  de  design:  um  espaço  de  endereço  para  dar  ao  processo  a  ilusão  de  ter  
sua  própria  memória  e  uma  thread  para  dar  ao  processo  a  ilusão  de  ter  sua  própria  CPU.  No  xv6,  um  processo  consiste  
em  um  espaço  de  endereço  e  uma  thread.  Em  sistemas  operacionais  reais,  um  processo  pode  ter  mais  de  uma  thread  
para  aproveitar  múltiplas  CPUs.
 2.6  Código:  iniciando  o  xv6,  o  primeiro  processo  e  chamada  do  sistema
 Para  tornar  o  xv6  mais  concreto,  descreveremos  como  o  kernel  inicia  e  executa  o  primeiro  processo.  Os  capítulos  
subsequentes  descreverão  os  mecanismos  que  aparecem  nesta  visão  geral  com  mais  detalhes.
 27
Machine Translated by Google
 Quando  o  computador  RISC-V  é  ligado,  ele  se  inicializa  e  executa  um  carregador  de  boot  armazenado  na  memória  
somente  leitura.  O  carregador  de  boot  carrega  o  kernel  xv6  na  memória.  Em  seguida,  no  modo  máquina,  a  CPU  executa  
o  xv6  a  partir  de  _entry  (kernel/entry.S:7).  O  RISC-V  começa  com  o  hardware  de  paginação  desabilitado:  endereços  
virtuais  são  mapeados  diretamente  para  endereços  físicos.
 O  carregador  carrega  o  kernel  xv6  na  memória  no  endereço  físico  0x80000000.  O  motivo  pelo  qual  ele  coloca  o  
kernel  em  0x80000000  em  vez  de  0x0  é  porque  o  intervalo  de  endereços  0x0:0x80000000  contém  dispositivos  de  E/S.
 As  instruções  em  _entry  configuram  uma  pilha  para  que  o  xv6  possa  executar  código  C.  O  xv6  declara  espaço  para  
uma  pilha  inicial,  stack0,  no  arquivo  start.c  (kernel/start.c:11).  O  código  em  _entry  carrega  o  registrador  de  ponteiro  de  
pilha  sp  com  o  endereço  stack0+4096,  o  topo  da  pilha,  porque  a  pilha  no  RISC-V  cresce  para  baixo.  Agora  que  o  kernel  
tem  uma  pilha,  _entry  chama  o  código  C  no  início  (kernel/start.c:21).
 A  função  start  realiza  algumas  configurações  permitidas  apenas  no  modo  máquina  e,  em  seguida,  alterna  para  o  
modo  supervisor.  Para  entrar  no  modo  supervisor,  o  RISC-V  fornece  a  instrução  mret.  Essa  instrução  é  mais  
frequentemente  usada  para  retornar  de  uma  chamada  anterior  do  modo  supervisor  para  o  modo  máquina.  start  não  
retorna  de  tal  chamada  e,  em  vez  disso,  configura  tudo  como  se  houvesse  uma:  define  o  modo  de  privilégio  anterior  como  
supervisor  no  registrador  mstatus,  define  o  endereço  de  retorno  como  main  escrevendo  o  endereço  de  main  no  registrador  
mepc,  desabilita  a  tradução  de  endereços  virtuais  no  modo  supervisor  escrevendo  0  no  registrador  de  tabela  de  páginas  
satp  e  delega  todas  as  interrupções  e  exceções  ao  modo  supervisor.
 Antes  de  entrar  no  modo  supervisor,  start  executa  mais  uma  tarefa:  programa  o  chip  do  relógio  
para  gerar  interrupções  de  tempo.  Com  essa  tarefa  concluída,  start  "retorna"  ao  modo  supervisor  
chamando  mret.  Isso  faz  com  que  o  contador  de  programa  mude  para  o  modo  principal  (kernel/main.c:11).
 Após  o  principal  (kernel/main.c:11)  inicializa  vários  dispositivos  e  subsistemas,  ele  cria  o  primeiro  processo  
chamando  userinit  (kernel/proc.c:233).  O  primeiro  processo  executa  um  pequeno  programa  escrito  em  assembly  RISC-V,  
que  faz  a  primeira  chamada  de  sistema  no  xv6.  initcode.S  (user/initcode.S:3)  carrega  o  número  para  a  chamada  do  
sistema  exec ,  SYS_EXEC  (kernel/syscall.h:8),  no  registro  a7  e,  em  seguida,  chama  ecall  para  entrar  novamente  no  
kernel.
 O  kernel  usa  o  número  no  registro  a7  em  syscall  (kernel/syscall.c:132)  para  chamar  a  chamada  de  sistema  desejada.  
A  tabela  de  chamadas  de  sistema  (kernel/syscall.c:107)  mapeia  SYS_EXEC  para  sys_exec,  que  o  kernel  invoca.  Como  
vimos  no  Capítulo  1,  exec  substitui  a  memória  e  os  registradores  do  processo  atual  por  um  novo  programa  (neste  caso, /
 init).
 Depois  que  o  kernel  conclui  o  exec,  ele  retorna  ao  espaço  do  usuário  no  processo /init .  init  (user/init.c:15)  Cria  um  
novo  arquivo  de  dispositivo  de  console,  se  necessário,  e  o  abre  como  descritores  de  arquivo  0,  1  e  2.  Em  seguida,  inicia  
um  shell  no  console.  O  sistema  está  funcionando.
 2.7  Modelo  de  Segurança
 Você  pode  se  perguntar  como  o  sistema  operacional  lida  com  códigos  maliciosos  ou  com  bugs.  Como  lidar  com  códigos  
maliciosos  é  mais  difícil  do  que  lidar  com  bugs  acidentais,  é  razoável  considerar  este  tópico  como  relacionado  à  
segurança.  Aqui  está  uma  visão  geral  das  premissas  e  objetivos  típicos  de  segurança  no  design  de  sistemas  operacionais.
 28
Machine Translated by Google
 O  sistema  operacional  deve  assumir  que  o  código  de  nível  de  usuário  de  um  processo  fará  o  possível  para  destruir  o  
kernel  ou  outros  processos.  O  código  do  usuário  pode  tentar  desreferenciar  ponteiros  fora  de  seu  espaço  de  endereço  
permitido;  pode  tentar  executar  quaisquer  instruções  RISC-V,  mesmo  aquelas  não  destinadas  ao  código  do  usuário;  pode  
tentar  ler  e  escrever  qualquer  registrador  de  controle  RISC-V;  pode  tentar  acessar  diretamente  o  hardware  do  dispositivo;  e  
pode  passar  valores  inteligentes  para  chamadas  de  sistema  na  tentativa  de  enganar  o  kernel  para  travar  ou  fazer  algo  
estúpido.  O  objetivo  do  kernel  é  restringir  cada  processo  do  usuário  para  que  tudo  o  que  ele  possa  fazer  seja  ler/escrever/
 executar  sua  própria  memória  de  usuário,  usar  os  32  registradores  RISC-V  de  uso  geral  e  afetar  o  kernel  e  outros  processos  
das  maneiras  que  as  chamadas  de  sistema  devem  permitir.  O  kernel  deve  impedir  quaisquer  outras  ações.  Isso  normalmente  
é  um  requisito  absoluto  no  projeto  do  kernel.
 As  expectativas  para  o  código  do  kernel  são  bem  diferentes.  O  código  do  kernel  é  assumido  como  escrito  por  
programadores  bem-intencionados  e  cuidadosos.  Espera-se  que  o  código  do  kernel  esteja  livre  de  bugs  e  certamente  não  
contenha  nada  malicioso.  Essa  suposição  afeta  a  forma  como  analisamos  o  código  do  kernel.  Por  exemplo,  existem  muitas  
funções  internas  do  kernel  (por  exemplo,  os  bloqueios  de  spin)  que  causariam  sérios  problemas  se  o  código  do  kernel  as  
utilizasse  incorretamente.  Ao  examinar  qualquer  trecho  específico  do  código  do  kernel,  queremos  nos  convencer  de  que  ele  
se  comporta  corretamente.  Assumimos,  no  entanto,  que  o  código  do  kernel  em  geral  está  escrito  corretamente  e  segue  
todas  as  regras  sobre  o  uso  das  funções  e  estruturas  de  dados  do  kernel.  No  nível  de  hardware,  presume-se  que  a  CPU,  a  
RAM,  o  disco  etc.  do  RISC-V  operem  conforme  anunciado  na  documentação,  sem  bugs  de  hardware.
 É  claro  que  na  vida  real  as  coisas  não  são  tão  simples.  É  difícil  evitar  que  um  código  de  usuário  inteligente  torne  um  
sistema  inutilizável  (ou  o  faça  entrar  em  pânico)  consumindo  recursos  protegidos  pelo  kernel  –  espaço  em  disco,  tempo  de  
CPU,  slots  na  tabela  de  processos,  etc.  Geralmente  é  impossível  escrever  código  livre  de  bugs  ou  projetar  hardware  livre  de  
bugs;  se  os  escritores  de  código  de  usuário  malicioso  estiverem  cientes  de  bugs  no  kernel  ou  no  hardware,  eles  os  
explorarão.  Mesmo  em  kernels  maduros  e  amplamente  utilizados,  como  o  Linux,  as  pessoas  descobrem  novas  
vulnerabilidades  continuamente  [1].  Vale  a  pena  projetar  salvaguardas  no  kernel  contra  a  possibilidade  de  ele  ter  bugs:  
asserções,  verificação  de  tipos,  páginas  de  proteção  de  pilha,  etc.  Finalmente,  a  distinção  entre  código  de  usuário  e  kernel  
às  vezes  é  confusa:  alguns  processos  privilegiados  em  nível  de  usuário  podem  fornecer  serviços  essenciais  e  efetivamente  
fazer  parte  do  sistema  operacional  e,  em  alguns  sistemas  operacionais,  o  código  de  usuário  privilegiado  pode  inserir  novo  
código  no  kernel  (como  com  os  módulos  carregáveis  do  kernel  do  Linux).
 2.8  Mundo  real
 A  maioria  dos  sistemas  operacionais  adotou  o  conceito  de  processo,  e  a  maioria  dos  processos  se  assemelha  aos  do  xv6.  
Os  sistemas  operacionais  modernos,  no  entanto,  suportam  várias  threads  dentro  de  um  processo,  permitindo  que  um  único  
processo  explore  múltiplas  CPUs.  Suportar  múltiplas  threads  em  um  processo  envolve  uma  série  de  mecanismos  que  o  xv6  
não  possui,  incluindo  possíveis  mudanças  na  interface  (por  exemplo,  o  clone  do  Linux,  uma  variante  do  fork)  para  controlar  
quais  aspectos  de  um  processo  as  threads  compartilham.
